% Options for packages loaded elsewhere
% Options for packages loaded elsewhere
\PassOptionsToPackage{unicode}{hyperref}
\PassOptionsToPackage{hyphens}{url}
\PassOptionsToPackage{dvipsnames,svgnames,x11names}{xcolor}
%
\documentclass[
  letterpaper,
  DIV=11,
  numbers=noendperiod]{scrartcl}
\usepackage{xcolor}
\usepackage{amsmath,amssymb}
\setcounter{secnumdepth}{-\maxdimen} % remove section numbering
\usepackage{iftex}
\ifPDFTeX
  \usepackage[T1]{fontenc}
  \usepackage[utf8]{inputenc}
  \usepackage{textcomp} % provide euro and other symbols
\else % if luatex or xetex
  \usepackage{unicode-math} % this also loads fontspec
  \defaultfontfeatures{Scale=MatchLowercase}
  \defaultfontfeatures[\rmfamily]{Ligatures=TeX,Scale=1}
\fi
\usepackage{lmodern}
\ifPDFTeX\else
  % xetex/luatex font selection
\fi
% Use upquote if available, for straight quotes in verbatim environments
\IfFileExists{upquote.sty}{\usepackage{upquote}}{}
\IfFileExists{microtype.sty}{% use microtype if available
  \usepackage[]{microtype}
  \UseMicrotypeSet[protrusion]{basicmath} % disable protrusion for tt fonts
}{}
\makeatletter
\@ifundefined{KOMAClassName}{% if non-KOMA class
  \IfFileExists{parskip.sty}{%
    \usepackage{parskip}
  }{% else
    \setlength{\parindent}{0pt}
    \setlength{\parskip}{6pt plus 2pt minus 1pt}}
}{% if KOMA class
  \KOMAoptions{parskip=half}}
\makeatother
% Make \paragraph and \subparagraph free-standing
\makeatletter
\ifx\paragraph\undefined\else
  \let\oldparagraph\paragraph
  \renewcommand{\paragraph}{
    \@ifstar
      \xxxParagraphStar
      \xxxParagraphNoStar
  }
  \newcommand{\xxxParagraphStar}[1]{\oldparagraph*{#1}\mbox{}}
  \newcommand{\xxxParagraphNoStar}[1]{\oldparagraph{#1}\mbox{}}
\fi
\ifx\subparagraph\undefined\else
  \let\oldsubparagraph\subparagraph
  \renewcommand{\subparagraph}{
    \@ifstar
      \xxxSubParagraphStar
      \xxxSubParagraphNoStar
  }
  \newcommand{\xxxSubParagraphStar}[1]{\oldsubparagraph*{#1}\mbox{}}
  \newcommand{\xxxSubParagraphNoStar}[1]{\oldsubparagraph{#1}\mbox{}}
\fi
\makeatother


\usepackage{longtable,booktabs,array}
\newcounter{none} % for unnumbered tables
\usepackage{calc} % for calculating minipage widths
% Correct order of tables after \paragraph or \subparagraph
\usepackage{etoolbox}
\makeatletter
\patchcmd\longtable{\par}{\if@noskipsec\mbox{}\fi\par}{}{}
\makeatother
% Allow footnotes in longtable head/foot
\IfFileExists{footnotehyper.sty}{\usepackage{footnotehyper}}{\usepackage{footnote}}
\makesavenoteenv{longtable}
\usepackage{graphicx}
\makeatletter
\newsavebox\pandoc@box
\newcommand*\pandocbounded[1]{% scales image to fit in text height/width
  \sbox\pandoc@box{#1}%
  \Gscale@div\@tempa{\textheight}{\dimexpr\ht\pandoc@box+\dp\pandoc@box\relax}%
  \Gscale@div\@tempb{\linewidth}{\wd\pandoc@box}%
  \ifdim\@tempb\p@<\@tempa\p@\let\@tempa\@tempb\fi% select the smaller of both
  \ifdim\@tempa\p@<\p@\scalebox{\@tempa}{\usebox\pandoc@box}%
  \else\usebox{\pandoc@box}%
  \fi%
}
% Set default figure placement to htbp
\def\fps@figure{htbp}
\makeatother





\setlength{\emergencystretch}{3em} % prevent overfull lines

\providecommand{\tightlist}{%
  \setlength{\itemsep}{0pt}\setlength{\parskip}{0pt}}



 


\usepackage{booktabs}
\usepackage{longtable}
\usepackage{array}
\usepackage{multirow}
\usepackage{wrapfig}
\usepackage{float}
\usepackage{colortbl}
\usepackage{pdflscape}
\usepackage{tabularx}
\usepackage{xltabular}
\usepackage{threeparttable}
\usepackage{threeparttablex}
\usepackage[normalem]{ulem}
\usepackage{makecell}
\usepackage{xcolor}
\KOMAoption{captions}{tableheading}
\makeatletter
\@ifpackageloaded{caption}{}{\usepackage{caption}}
\AtBeginDocument{%
\ifdefined\contentsname
  \renewcommand*\contentsname{Table of contents}
\else
  \newcommand\contentsname{Table of contents}
\fi
\ifdefined\listfigurename
  \renewcommand*\listfigurename{List of Figures}
\else
  \newcommand\listfigurename{List of Figures}
\fi
\ifdefined\listtablename
  \renewcommand*\listtablename{List of Tables}
\else
  \newcommand\listtablename{List of Tables}
\fi
\ifdefined\figurename
  \renewcommand*\figurename{Figure}
\else
  \newcommand\figurename{Figure}
\fi
\ifdefined\tablename
  \renewcommand*\tablename{Table}
\else
  \newcommand\tablename{Table}
\fi
}
\@ifpackageloaded{float}{}{\usepackage{float}}
\floatstyle{ruled}
\@ifundefined{c@chapter}{\newfloat{codelisting}{h}{lop}}{\newfloat{codelisting}{h}{lop}[chapter]}
\floatname{codelisting}{Listing}
\newcommand*\listoflistings{\listof{codelisting}{List of Listings}}
\makeatother
\makeatletter
\makeatother
\makeatletter
\@ifpackageloaded{caption}{}{\usepackage{caption}}
\@ifpackageloaded{subcaption}{}{\usepackage{subcaption}}
\makeatother
\usepackage{bookmark}
\IfFileExists{xurl.sty}{\usepackage{xurl}}{} % add URL line breaks if available
\urlstyle{same}
% fallback for those not using the hyperref driver hyperxmp:
\makeatletter
\@ifundefined{xmpquote}{\newcommand{\xmpquote}[1]{#1}}{}
\makeatother
\hypersetup{
  pdftitle={PERSONALITY PSYCHOLOGY LAB},
  colorlinks=true,
  linkcolor={blue},
  filecolor={Maroon},
  citecolor={Blue},
  urlcolor={Blue},
  pdfcreator={LaTeX via pandoc}}


\title{PERSONALITY PSYCHOLOGY LAB}
\usepackage{etoolbox}
\makeatletter
\providecommand{\subtitle}[1]{% add subtitle to \maketitle
  \apptocmd{\@title}{\par {\large #1 \par}}{}{}
}
\makeatother
\subtitle{PSYC BC2124\\
Barnard College\\
Department of Psychology\\
Fall 2026}
\author{}
\date{}
\begin{document}
\maketitle


\textbf{Pre-requisites}: PSYC BC1001 Introduction to Psychology; PSYC
BC1101 Statistics; PSYC BC1020 Research Methods\\
\textbf{Co-requisite}: PSYC BC2125 Personality Psychology

\subsubsection{Time, venue \& instructor}\label{time-venue-instructor}

Tuesday 9:10AM-12:00PM, Milbank 410

Instructor: Dr.~Annie Vu
(\href{mailto:avu@barnard.edu}{\nolinkurl{avu@barnard.edu}})\\
Office hour: MW: 1:00pm-2:00pm; T: 1:00pm-3:00pm. I am happy to meet in person or on Zoom by appointment (request via email or in person).



\subsubsection{Overview}\label{overview}

This lab will usually be taken concurrently with BC2125 Personality
lecture. It will expand upon some of the theoretical, methodological and
analytic issues introduced there, as well as giving you the opportunity
to explore topics from within and beyond those covered in the lectures
in greater depth through hands-on experience with research design and
data analysis.

The semester is broken into 3 projects, each illustrating a common
methodological and analytic approach in personality psychology. Through
these projects you will gain experience in formulating psychological
research questions relating to personality; analyzing and visualizing
data; and interpreting and communicating your findings. Projects will be
undertaken in small groups; group members will collaborate on design and
analysis. Each project will culminate in an in-class group presentation
and submission of a brief individual write up of your project.

\subsubsection{Schedule and deadlines}\label{schedule-and-deadlines}

\begin{xltabular}{\linewidth}{>{\raggedright\arraybackslash}X>{\raggedright\arraybackslash}X}
\toprule
Date & Topic\\
\midrule
9/8 & Course overview\\
\addlinespace[0.3em]
\multicolumn{2}{l}{\textbf{Project 1: t-test}}\\
9/15 & Project planning\\
9/22 & Data cleaning\\
9/29 & Analysis\\
10/6 & Presentations\\
\addlinespace[0.3em]
\multicolumn{2}{l}{\textbf{Project 2: Correlation}}\\
10/13 & Project planning\\
10/20 & Data cleaning \& analysis\\
10/27 & Visualization \& reporting\\
\textcolor[HTML]{BBBBBB}{\em{11/4}} & \textcolor[HTML]{BBBBBB}{\em{No class (academic holiday)}}\\
11/10 & Presentations\\
\addlinespace[0.3em]
\multicolumn{2}{l}{\textbf{Project 3: Regression}}\\
11/17 & Project planning\\
11/24 & Data cleaning \& analysis\\
12/1 & Visualization \& reporting\\
12/8 & Presentations\\
\bottomrule
\end{xltabular}

Written project reports are due one week following the in-class
presentation (before the next week's lab session) on the following
dates.

{\def\LTcaptype{none} % do not increment counter
\begin{longtable}[]{@{}ll@{}}
\toprule\noalign{}
& Due Date \\
\midrule\noalign{}
\endhead
\bottomrule\noalign{}
\endlastfoot
Project 1 report & 10/13 \\
Project 2 report & 11/17 \\
Project 3 report & 12/15 \\
\end{longtable}
}

\subsubsection{Class format \&
participation}\label{class-format-participation}

Labs are substantially more interactive and discussion-based than the
traditional lecture format, and depend on everyone's active
participation in class discussions and activity as well as group work
focused around the projects. Your active participation across the
semester will therefore contribute a substantial portion of your grade.

If you have questions, thoughts, or ideas you want to share, feel free
to do so at any time (while keeping within the bounds of polite
conversation, obviously--don't interrupt or talk over other people! But
do feel free to respond to others without having to raise your hand or
wait to be called on). Everyone will get the most out of this lab when
the discussion can develop organically and everyone feels free to be
part of the conversation if \& when they have something to add.

Being part of the in-person discussion is one obvious way to
participate, but it's not the only way. Different people have different
styles of participation, and the lab is designed to try and accommodate
and encourage different approaches. Your level of engagement with your
project partner(s), your TA and Prof.~Brotherton as you work through
your projects is also an important form of participation. You can also
participate by coming to office hours.

At a minimum (i.e.~for a passing grade), I'll be looking for some form
of participation (loosely defined) from you every week. Higher
participation grades will be earned through regular, enthusiastic,
productive participation (note that quality is more important than
quantity).

\subsubsection{Workload}\label{workload}

As a general rule for the amount of time students should expect to
commit to classes, the college suggests three hours per week in or
outside of class per credit. Since this class is worth 2.5 credits, that
corresponds to 7.5 hours per week, split between time in the classroom
and time spent completing the associated assignments.

\subsubsection{Final Grades}\label{final-grades}

Your numeric score for the course is a product of your scores for each
assignment, weighted as follows:

{\def\LTcaptype{none} % do not increment counter
\begin{longtable}[]{@{}lr@{}}
\toprule\noalign{}
& Weight (\%) \\
\midrule\noalign{}
\endhead
\bottomrule\noalign{}
\endlastfoot
Participation (over the course of the semester) & 10 \\
Presentations & 30 \\
Reports & 60 \\
\end{longtable}
}

Final grades are determined according to the following boundaries:

\begin{verbatim}
Letter grade:  A+ A  A- B+ B  B- C+ C  C- D  F
Numeric score: 97 93 90 87 83 80 77 73 70 60 <60
\end{verbatim}

\subsection{Course policies}\label{course-policies}

\subsubsection{Attendance \& timeliness}\label{attendance-timeliness}

In-person attendance of every lab session is expected, and you should
expect to stay for the full duration of the lab. Normally it is
departmental policy to remove students who miss more than two lab
sessions from the course; however, given ongoing revisions to
college-wide health-related policies, exceptions may be made. If you are
feeling unwell, you should not come to class and notify me of
nonattendance before class if possible.

When you are attending, please arrive on time for class. Frequent
lateness will impact your participation grade.

\subsubsection{Assignment late policy}\label{sec-late-policy}

Presentations will be given in-class on the dates listed in the class
schedule. It is expected that all group members contribute equally to
the presentation preparation and delivery. Therefore, all group members
will receive the same group score, except in exceptional circumstances
when it is clear that contributions were not equal. If a group member is
unable to present in-person on the day of the presentation due to
unforeseen circumstances, please notify Prof.~Brotherton. The
presentation will be given by the remaining group members, and the
absent member may be assigned the group score, assuming equal
contributions to the preparation of the presentation.

Written project reports are due one week following the in-class
presentation, i.e.~before the next week's lab session--see dates below.
A grade penalty of 5 points will be applied for each day (or part
thereof) that an assignment is late (up to a maximum of 6 days; work not
submitted before the next lab will receive a score of zero). For
example, if your work is A+ quality but is submitted a day-and-a-half
late, you will only receive a B+. This policy is intended to incentivize
timely submission while easing the stress of genuine emergencies. When
things come up that prevent timely submission you can prioritize
accordingly, knowing that a small penalty on one assignment for this lab
will not tank your final grade.

\subsubsection{Academic integrity}\label{academic-integrity}

Students are expected to follow the Barnard Honor Code, available at
\url{https://barnard.edu/honor-code}.

Note that while you will collaborate with group members on the design,
analysis, and presentation of research projects, you may not collaborate
on the written report: each group member must write their own individual
reports.

\subsubsection{Use of generative AI
(e.g.~ChatGPT)}\label{use-of-generative-ai-e.g.-chatgpt}

Generative AI is an extremely useful tool for many purposes; real
researchers are of course using it to help with their research. For this
class, you are welcome to use AI tools to help with idea generation,
help you understand readings, help formulate your ideas for written
reports, etc. However, note the emphasis on \emph{helping} rather than
doing things for you. Passing someone else's work off as your own is
plagiarism and a violation of the honor code, and that includes the work
of AI assistants. AI tools can enhance, but should not substitute for,
your own efforts. Please be aware that it is generally obvious when
someone has put little thought of their own into their work. Relying on
ChatGPT (or similar) to do most of the work for you will result in bad
work (and will accordingly earn you a bad grade).

\subsubsection{Academic accommodations and general
wellness}\label{academic-accommodations-and-general-wellness}

It is always important to recognize the different pressures, burdens,
and stressors you may be facing, whether personal, emotional, physical,
financial, mental, or academic. The faculty and administration recognize
this, and are prepared to provide assistance to students in need. I
encourage you to seek advice from your advisor, Dean, the
\href{https://barnard.edu/disability-services}{Center for Accessibility
Resources \& Disability Services (CARDS)}, or
\href{https://barnard.edu/health-wellness}{Barnard Health \& Wellness}
as needed. Please let me know of any issues you wish to share with me
that you feel are impacting your ability to complete the course to the
best of your ability. Though it isn't always easy, it is better to
proactively seek help rather than letting problems build up.




\end{document}
